\documentclass[10pt]{article}

\usepackage{amsmath}
%
\usepackage{amssymb}
%
\usepackage{enumerate}
%
\usepackage[width=7in,height=10in,top=0.75in,centering]{geometry}
%
\usepackage[dvipsnames]{xcolor}
%
\usepackage{tikz}
%
\usepackage{multicol}
%
\usepackage{hyperref}
%
\usepackage{marginnote}
%
%%%%%%%%%%%%%%%% HEADER %%%%%%%%%%%%%%%%%%%%%
\usepackage{fancyhdr}
\pagestyle{fancy}
%%
%\lfoot{Left footer}
%\cfoot{Center of footer}
%\rfoot{Right of footer}
%%
\lhead{Survey of Calculus I}
\chead{Week 3 Workshop}
\rhead{\thepage}
%
%%%%%%%%%%%%%%%% HEADER %%%%%%%%%%%%%%%%%%%%%
%
%%%%%%%%%%%%%%%% PGFPLOTS %%%%%%%%%%%%%%%%%%%
\usepackage{pgfplots}
%
\pgfplotsset{every axis/.append style={
                    axis x line=middle,    % put the x axis in the middle
                    axis y line=middle,    % put the y axis in the middle
                    axis line style={<->,color=black}, % arrows on the axis
                    %xlabel={$x$},          % default put x on x-axis
                    ylabel={$y$},          % default put y on y-axis
            }}
%
 \pgfplotsset{compat=1.14}
%
 \usepgfplotslibrary{fillbetween}
%%%%%%%%%%%%%%%% PGFPLOTS %%%%%%%%%%%%%%%%%%%
%
\parindent0pt
\fboxsep=0.25cm
%
\newcommand{\Heading}[2]{\fbox{{\Large\textbf{#1}}}\hfill\fbox{\textbf{#2}}\par\vspace{0.1in}}
%
\newcommand{\Name}[1]{\textbf{Name}:\quad#1\hrulefill}
%
\newcommand{\Target}[1]{\textbf{Learning Target(s):}\quad\fbox{#1}\par\medskip}
%
\DeclareMathOperator{\arcsec}{arcsec}
\DeclareMathOperator{\arccot}{arccot}
\DeclareMathOperator{\arccsc}{arccsc}
%
\newcounter{example}[section]
\newcounter{exercise}[section]
%
\newcommand{\important}[2]{\begin{center}\fbox{\begin{minipage}{.97\textwidth}\noindent\textbf{#1}\par#2\end{minipage}}\end{center}}
%
\newcommand{\defn}[1]{\begin{center}\fbox{\begin{minipage}{.97\textwidth}\reversemarginpar\marginnote{\textbf{Definition}}#1\end{minipage}}\end{center}}
%
\newcommand{\example}[1]{\vspace{0.1in}{\textcolor{blue}{\textbf{Example \stepcounter{example}\theexample.}}}\hspace{0.1in}\reversemarginpar\marginnote{\fbox{\bf #1}}}
%
\newcommand{\exercise}[1]{\vspace{0.1in}{\textcolor{blue}{\textbf{Exercise \stepcounter{exercise}\Alph{exercise}.}}}\hspace{0.1in}\reversemarginpar\marginnote{\fbox{\bf #1}}}

%%%%%%%%%%%% BEGIN DOCUMENT %%%%%%%%%%%%%%%
%%%%%%%%%%%% BEGIN DOCUMENT %%%%%%%%%%%%%%%
%%%%%%%%%%%% BEGIN DOCUMENT %%%%%%%%%%%%%%%
%%%%%%%%%%%% BEGIN DOCUMENT %%%%%%%%%%%%%%%
%%%%%%%%%%%% BEGIN DOCUMENT %%%%%%%%%%%%%%%

\begin{document}

\thispagestyle{empty}

\Heading{MATH 1190 \& 1210 Survey of Calculus I}{\begin{minipage}{0.25\textwidth}\begin{center} Week 3 Workshop:\\ Simplifying Exponents\\ Limits \& Forms\end{center}\end{minipage}}

\vspace{0.1in}

\Name{} \quad \Target{{\bf\color{RoyalBlue}L3}}

\hrulefill\\

%In lesson 4, we computed limits involving expressions with exponents -- such as $x^{-1/2}$ and $e^{2x}$ -- and in lesson 6 we will work with these kinds of functions again. So, let's get some practice rewriting expressions with exponents!

%\begin{enumerate}
%\item[1.] Rewrite the following expressions with only positive exponents. Combine exponential expressions using exponent rules, when possible.\\
%\end{enumerate}

%\begin{align*}
%x^{-7} &=\rule[-8pt]{1in}{0.4pt} & \frac12x^{-3} &=\rule[-8pt]{1in}{0.4pt} & \dfrac{1}{4x^{-2}} &=\rule[-8pt]{1in}{0.4pt}\\[0.75in]
%
% \dfrac{2x^{-2}}{5x} &=\rule[-8pt]{1in}{0.4pt} & \dfrac{2x^{-4}}{3x^{-5}} &=\rule[-8pt]{1in}{0.4pt} & \dfrac{4x^2}{7x^{-6}} &=\rule[-8pt]{1in}{0.4pt}\\[0.75in]
% %
% \frac23e^{-x} &=\rule[-8pt]{1in}{0.4pt} & \dfrac{e^{2x}}{2e^{-x}} &=\rule[-8pt]{1in}{0.4pt} & \dfrac{4e^{-x}}{7e^x} &=\rule[-8pt]{1in}{0.4pt}\\[0.75in]
% \end{align*}

% \begin{enumerate}
% \item[2.] Rewrite the following expressions using exponent notation. For example, $\dfrac{4}{\sqrt[3]{x^2}}$ would be rewritten as $4x^{-2/3}$. Combine exponential expressions using exponent rules, when possible, and write your final answer using only positive exponents.\\
% \end{enumerate}

% \begin{align*}
% \frac{\sqrt[4]{x^7}}{2} &=\rule[-8pt]{1in}{0.4pt} & \frac{1}{2\sqrt{x^3}} &=\rule[-8pt]{1in}{0.4pt} & \dfrac{3x^{5/3}}{2\sqrt[3]{x}} &=\rule[-8pt]{1in}{0.4pt}\\[0.75in]
% %
% \dfrac{\sqrt[3]{x+1}}{(x+1)^{-2}} &=\rule[-8pt]{1in}{0.4pt} & \dfrac{x+1}{\sqrt{x^2+1}} &=\rule[-8pt]{1in}{0.4pt} & \dfrac{e^{-x}}{\sqrt{e^x}} &=\rule[-8pt]{1in}{0.4pt}
% \end{align*}

% %%%%%%%%%
% \newpage%
% %%%%%%%%%

% \begin{enumerate}
% \item[3.] Evaluate the following limits. Show your work, using forms when necessary.
% \end{enumerate}

% \begin{align*}
% \lim_{x \rightarrow 5} 5x^2 &=\rule[-8pt]{1in}{0.4pt} & \lim_{x \rightarrow \infty} 5x^2 &=\rule[-8pt]{1in}{0.4pt} & \lim_{x \rightarrow \infty} \frac{3x+5}{e^x} &=\rule[-8pt]{1in}{0.4pt}\\[1in]
% %
% \lim_{x \rightarrow \infty} \frac{e^x}{3x+5} &=\rule[-8pt]{1in}{0.4pt} & \lim_{x \rightarrow \infty} \frac{x^2}{x^5} &=\rule[-8pt]{1in}{0.4pt} & \lim_{x \rightarrow 5} \frac{3x^2 + 8}{x^9}&=\rule[-8pt]{1in}{0.4pt}\\[1in]
% \end{align*}

\begin{enumerate}
\item[1.] Evaluate the following limits. Show your work, using forms when necessary.
\end{enumerate}

\begin{align*}
\lim_{x \rightarrow 4} \sqrt{x} - 5 &=\rule[-8pt]{1in}{0.4pt} & \lim_{x \rightarrow \infty} \sqrt{x} - 5 &=\rule[-8pt]{1in}{0.4pt} & \lim_{x \rightarrow \infty} (1-x)^{-3} &=\rule[-8pt]{1in}{0.4pt}\\[1in]
%
\lim_{x \rightarrow 3} (3-x)^{-3} &=\rule[-8pt]{1in}{0.4pt} & \lim_{x \rightarrow \infty} \frac{x^2 + 3x^4}{x^5} &=\rule[-8pt]{1in}{0.4pt} & \lim_{x \rightarrow \infty} x - x + 100 &=\rule[-8pt]{1in}{0.4pt}\\[1in]
\end{align*} 




\begin{enumerate}
\item[2.]\reversemarginpar\marginnote{\fbox{\bf\color{RoyalBlue}L3}} For each of the limits below, first write the form, then evaluate the limits. Show your work.
    \begin{enumerate}
    \item[(a)] $\displaystyle\lim_{x\to\infty}\frac{x^2+\sqrt{x}-1}{e^x-1}$ \hfill {\bf form:}\rule[-8pt]{0.5in}{0.4pt} \hfill {\bf value:}\rule[-8pt]{0.5in}{0.4pt}
        \vfill

    \item[(b)] $\displaystyle\lim_{x\to0}\frac{x^2+\sqrt{x}-1}{e^x-1}$ \hfill {\bf form:}\rule[-8pt]{0.5in}{0.4pt} \hfill {\bf value:}\rule[-8pt]{0.5in}{0.4pt}
        \vfill
    \end{enumerate}
\end{enumerate}

\begin{enumerate}
\item[3.]\reversemarginpar\marginnote{\fbox{\bf\color{RoyalBlue}L3}} Give a function satisfying each set of criteria
    \begin{enumerate}
    \item[(a)] As $x$ approaches infinity, $f(x)$ is of the form $\frac \infty \infty$ and $\lim_{x\to \infty}f(x) = 3$
        \vfill

    \item[(b)] As $x$ approaches negative infinity, $f(x)$ is of the form $\frac \infty \infty$ and $\lim_{x\to -\infty}f(x) = \infty$ 
        \vfill

    \item[(c)] Challenge problem: $\lim_{x\to -\infty}f(x) = 3$ and $\lim_{x\to \infty}f(x) = 6$ 
        \vfill
    \end{enumerate}
\end{enumerate}

\end{document}